%%%%%%%% ICML 2026 EXAMPLE LATEX SUBMISSION FILE %%%%%%%%%%%%%%%%%

\documentclass{article}

% Recommended, but optional, packages for figures and better typesetting:
\usepackage{microtype}
\usepackage{graphicx}
\usepackage{subcaption}
\usepackage{booktabs} % for professional tables

% hyperref makes hyperlinks in the resulting PDF.
% If your build breaks (sometimes temporarily if a hyperlink spans a page)
% please comment out the following usepackage line and replace
% \usepackage{icml2026} with \usepackage[nohyperref]{icml2026} above.
\usepackage{hyperref}


% Attempt to make hyperref and algorithmic work together better:
\newcommand{\theHalgorithm}{\arabic{algorithm}}

% Use the following line for the initial blind version submitted for review:
\usepackage{icml2026}

% For preprint, use
% \usepackage[preprint]{icml2026}

% If accepted, instead use the following line for the camera-ready submission:
% \usepackage[accepted]{icml2026}

\usepackage{amsmath}
\usepackage{amssymb}
\usepackage{mathtools}
\usepackage{amsthm}


% if you use cleveref..
\usepackage[capitalize,noabbrev]{cleveref}

%%%%%%%%%%%%%%%%%%%%%%%%%%%%%%%%
% THEOREMS
%%%%%%%%%%%%%%%%%%%%%%%%%%%%%%%%
\theoremstyle{plain}
\newtheorem{theorem}{Theorem}[section]
\newtheorem{proposition}[theorem]{Proposition}
\newtheorem{lemma}[theorem]{Lemma}
\newtheorem{corollary}[theorem]{Corollary}
\theoremstyle{definition}
\newtheorem{definition}[theorem]{Definition}
\newtheorem{assumption}[theorem]{Assumption}
\theoremstyle{remark}
\newtheorem{remark}[theorem]{Remark}

% Todonotes is useful during development; simply uncomment the next line
%    and comment out the line below the next line to turn off comments
%\usepackage[disable,textsize=tiny]{todonotes}
\usepackage[textsize=tiny]{todonotes}

% The \icmltitle you define below is probably too long as a header.
% Therefore, a short form for the running title is supplied here:
\icmltitlerunning{Sparsity-Preserving Joint Alignment}

\begin{document}

\twocolumn[
  \icmltitle{Sparsity-Preserving Joint Alignment: Overcoming Activation Variance \\
    Collapse and Representation Drift in Multi-Task Model Merging}

  % It is OKAY to include author information, even for blind submissions: the
  % style file will automatically remove it for you unless you've provided
  % the [accepted] option to the icml2026 package.

  % List of affiliations: The first argument should be a (short) identifier you
  % will use later to specify author affiliations Academic affiliations
  % should list Department, University, City, Region, Country Industry
  % affiliations should list Company, City, Region, Country

  % You can specify symbols, otherwise they are numbered in order. Ideally, you
  % should not use this facility. Affiliations will be numbered in order of
  % appearance and this is the preferred way.
  \icmlsetsymbol{equal}{*}

  \begin{icmlauthorlist}
    \icmlauthor{Firstname1 Lastname1}{equal,yyy}
    \icmlauthor{Firstname2 Lastname2}{equal,yyy,comp}
    \icmlauthor{Firstname3 Lastname3}{comp}
    \icmlauthor{Firstname4 Lastname4}{sch}
    \icmlauthor{Firstname5 Lastname5}{yyy}
    \icmlauthor{Firstname6 Lastname6}{sch,yyy,comp}
    \icmlauthor{Firstname7 Lastname7}{comp}
    %\icmlauthor{}{sch}
    \icmlauthor{Firstname8 Lastname8}{sch}
    \icmlauthor{Firstname8 Lastname8}{yyy,comp}
    %\icmlauthor{}{sch}
    %\icmlauthor{}{sch}
  \end{icmlauthorlist}

  \icmlaffiliation{yyy}{Department of XXX, University of YYY, Location, Country}
  \icmlaffiliation{comp}{Company Name, Location, Country}
  \icmlaffiliation{sch}{School of ZZZ, Institute of WWW, Location, Country}

  \icmlcorrespondingauthor{Firstname1 Lastname1}{first1.last1@xxx.edu}
  \icmlcorrespondingauthor{Firstname2 Lastname2}{first2.last2@www.uk}

  % You may provide any keywords that you find helpful for describing your
  % paper; these are used to populate the "keywords" metadata in the PDF but
  % will not be shown in the document
  \icmlkeywords{Machine Learning, ICML}

  \vskip 0.3in
]

% this must go after the closing bracket ] following \twocolumn[ ...

% This command actually creates the footnote in the first column listing the
% affiliations and the copyright notice. The command takes one argument, which
% is text to display at the start of the footnote. The \icmlEqualContribution
% command is standard text for equal contribution. Remove it (just {}) if you
% do not need this facility.

% Use ONE of the following lines. DO NOT remove the command.
% If you have no special notice, KEEP empty braces:
\printAffiliationsAndNotice{}  % no special notice (required even if empty)
% Or, if applicable, use the standard equal contribution text:
% \printAffiliationsAndNotice{\icmlEqualContribution}

\begin{abstract}
  Multi-task model merging has emerged as a powerful, training-free paradigm to combine specialized expert models into a single unified backbone. However, direct parameter averaging typically triggers activation variance collapse, while task conflicts distort deep feature representations. Existing solutions either calibrate backbone statistics (which often disrupts sparse activation pathways like those in ReLU networks) or adapt classification heads on tiny calibration sets (which fails to recover backbone variance decay). To resolve these bottlenecks, we propose Sparsity-Preserving Joint Alignment (SPJA), a unified framework that combines Native Task-Agnostic Activation Calibration (N-TAAC) with decision-boundary head adaptation. SPJA leverages N-TAAC to update the backbone's BatchNorm statistics on a joint calibration set, restoring representational flow globally and task-agnostically while naturally preserving ReLU sparsity (since BatchNorm is applied prior to ReLU activations). Concurrently, SPJA adapts classification heads to align final decision boundaries in this calibrated representational space, avoiding task-conditional hooks or routing in the backbone during inference. Through extensive, multi-seed empirical sweeps across diverse computer vision tasks (MNIST, Fashion-MNIST, CIFAR-10), we demonstrate that SPJA consistently outperforms existing calibration or head-adaptation baselines, establishing a robust new standard for representational alignment in post-hoc multi-task model merging.
\end{abstract}

\section{Introduction}
Deploying multi-task neural network architectures is a central challenge in modern deep learning \cite{caruana1997multitask, ruder2017overview, misra2016cross, ilharco2022editing}. While training a unified model from scratch is computationally prohibitive, post-hoc model merging has emerged as a powerful, training-free alternative \cite{wortsman2022model, yadav2023ties, yang2024survey, khan2024deep}. By directly interpolating the weights of specialized downstream "expert" models fine-tuned independently from a shared pre-trained base model, model merging combines distinct capabilities with zero training cost \cite{matena2021merging, jin2022regmean, ainsworth2022git, stoica2023zipit}.

Despite its elegance, model merging suffers from parameter interference. When expert updates are averaged linearly (as in Weight Averaging) or linearly combined (as in Task Arithmetic), task-specific parameter structures collide \cite{izmailov2018averaging, garipov2018loss, draxler2018essentially}. In activation space, this parameter conflict triggers an exponential decay of activation variance in deeper layers, a phenomenon known as activation variance collapse \cite{jordan2022repair}. This collapse restricts deep features to a near-zero variance shell, degrading multi-task performance to near-random guessing.

To mitigate this, prior work has proposed representation calibration and activation restoration, notably REPAIR \cite{jordan2022repair} and Task-Conditional Activation Calibration (TCAC). These methods apply channel-wise affine transformations using statistics estimated over a small calibration set. However, channel-wise calibration has serious drawbacks: (i) in ReLU-based architectures, sparse channels with zero activation on small calibration sets cause division-by-zero or scale noise (the "sparsity trap"), and (ii) subtracting the mean shift alters the sign of non-negative activations, violating ReLU sparsity and disrupting expert routing pathways. Layer-wise Scaling-only Calibration (LSC) resolves this by utilizing a single global scalar per layer, stabilizing variance while perfectly preserving ReLU sparsity.

Separately, recent representational analysis reveals that backbone features remain highly intact in shallow and middle layers, with representation drift and task conflicts primarily localizing in the final layers and classification heads. Consequently, adapting only the classification heads (Head-only SFT or TTA) on a tiny subset of 128 samples is highly data-efficient and recovers the vast majority of performance, suggesting that intermediate activation calibration may be over-engineered.

In this work, we deconstruct both paradigms and ask: \textit{Are representation calibration and decision-boundary head adaptation mutually exclusive, or are they complementary?} We hypothesize that they address orthogonal issues: while backbone calibration heals global signal propagation (variance collapse), classification head adaptation corrects localized decision-boundary misalignment. By combining them under a mathematically consistent framework, we can achieve outstanding performance recovery that completely outperforms either approach individually.

To validate this, we propose \textbf{Sparsity-Preserving Joint Alignment (SPJA)}. SPJA first calibrates backbone signal propagation using Native Task-Agnostic Activation Calibration (N-TAAC), which performs a single forward pass over a joint calibration set in native training mode with frozen weights, updating the BatchNorm running statistics. This restores backbone representation space globally and task-agnostically while naturally preserving ReLU sparsity (since BatchNorm is applied prior to ReLU activations). Concurrently, SPJA adapts task heads (via Supervised SFT, SPJA-SFT, or Unsupervised TTA, SPJA-TTA) to align decision boundaries with this joint, task-agnostic calibrated representation space.

We conduct an exhaustive evaluation of SPJA alongside all relevant baselines (NONE, N-TAAC, LSC, Head-SFT, and Head-TTA) on a diverse vision benchmark (MNIST, Fashion-MNIST, and CIFAR-10). Our results across multiple calibration sizes ($N \in \{32, 128, 512\}$), Task Arithmetic coefficients, and independent random seeds provide overwhelming empirical evidence for our claims, establishing SPJA as the new state-of-the-art standard for multi-task representation alignment.

\section{Related Work}
\textbf{Parameter-Space Model Merging:} Model merging and parameter interpolation originate from model optimization and federated learning, such as Stochastic Weight Averaging \cite{izmailov2018averaging}, Model Soups \cite{wortsman2022model}, and federated aggregation schemes (FedAvg \cite{fedavg2017}, FedProx \cite{li2020federated}, FedOpt \cite{reddi2020adaptive}). For combining fine-tuned experts, Weight Averaging (WA) and Task Arithmetic (TA) \cite{ilharco2022editing} are standard baselines. To combat parameter interference, advanced methods like TIES-Merging \cite{yadav2023ties} resolve sign conflicts and prune small updates, DARE \cite{yu2024language} applies random sparsification, Git Re-Basin \cite{ainsworth2022git} permutes units to align basins, and ZipIt! \cite{stoica2023zipit} merges disparate topologies. Comprehensive surveys \cite{yang2024survey, khan2024deep, wang2024scaling} map this rapidly expanding space. While cheap, these parameter-space methods ignore the downstream activation distortions caused by parameter-level merging.

\textbf{Activation Calibration:} To heal representational distortions, REPAIR \cite{jordan2022repair} rescales permuted activations to restore statistics, and TCAC extends this using task-conditional scaling. Modern activation-aware merging techniques like Chain of Merges \cite{buzzega2025chain} and Neural Merger \cite{horoi2024neural} align features on small calibration sets, while foldable supernets \cite{kinderman2024foldable} map activations to a shared workspace. Similar ideas are used in quantization calibration, such as SmoothQuant \cite{xiao2023smoothquant} and AWQ \cite{lin2023awq}. To eliminate the complex hooks and task-ID dependencies of TCAC, Native Task-Agnostic Activation Calibration (N-TAAC) uses training mode to update BatchNorm statistics. However, channel-wise calibration is highly sensitive to the sparsity trap of ReLU networks. LSC resolves this by using a single positive standard deviation ratio globally per layer, preserving non-negativity and sparsity.

\textbf{Head-Only Adaptation:} Head adaptation baselines originate from parameter-efficient fine-tuning (PEFT) \cite{houlsby2019parameter, hu2021lora, li2021prefix, lester2021power, zaken2021bitfit} and test-time adaptation (TTA) \cite{wang2021tent, sun2020test, wang2022continual, boudiaf2022parameter}, which often leverage knowledge distillation principles \cite{hinton2015distilling} to transfer learned representations. Recent deconstructive studies show that backbone features of merged models are mostly preserved (CKA > 0.95), with representation drift localizing primarily in the final layers. Therefore, adapting only the classification heads on tiny calibration subsets recovers almost all performance, avoiding intrusive hooks or test-time coefficient optimization.

\section{Methodology}
\subsection{Problem Setup}
Let $\Theta_{\text{pre}}$ denote the weights of a pre-trained base model. We are given $K$ expert models, independently fine-tuned from $\Theta_{\text{pre}}$ on separate downstream tasks $k \in \{1, \dots, K\}$, resulting in task-specific weights $\Theta_1, \dots, \Theta_K$. Each expert model consists of a shared backbone $\theta_k$ and a task-specific classification head $h_k$. 

Our goal is to construct a unified multi-task model with a single merged backbone $\theta_{\text{merged}}$ and the original task-specific heads $h_1, \dots, h_K$. During inference on a test sample $x$ belonging to task $k$, the model uses the merged backbone and task-specific head $h_k$ to make predictions: $f_k(x) = h_k(\Phi(x; \theta_{\text{merged}}))$.

\subsection{Backbone Merging Baselines}
We consider two standard backbone merging baselines:
\begin{enumerate}
    \item \textbf{Weight Averaging (WA):} The merged backbone is computed as the simple element-wise arithmetic mean:
    \begin{equation}
        \theta_{\text{merged}}^{\text{WA}} = \frac{1}{K}\sum_{k=1}^K \theta_k
    \end{equation}
    \item \textbf{Task Arithmetic (TA):} We extract task vectors $\tau_k = \theta_k - \theta_{\text{pre}}$, and the merged weights are:
    \begin{equation}
        \theta_{\text{merged}}^{\text{TA}} = \theta_{\text{pre}} + \lambda \sum_{k=1}^K \tau_k
    \end{equation}
    where $\lambda \in [0, 1]$ is a global scaling coefficient.
\end{enumerate}

\subsection{Sparsity-Preserving Joint Alignment (SPJA)}
Our proposed framework, Sparsity-Preserving Joint Alignment (SPJA), consists of two core steps: backbone calibration via Native Task-Agnostic Activation Calibration (N-TAAC), followed by N-TAAC-guided classification head adaptation.

\subsubsection{Step 1: Joint Backbone Calibration via N-TAAC}
Instead of task-conditional scaling (which requires task-ID routing in the backbone during inference), SPJA leverages Native Task-Agnostic Activation Calibration (N-TAAC) to calibrate the merged model's backbone activations globally. Let $D_{\text{cal}} = \cup_{k=1}^K D_{\text{cal}}^{(k)}$ be the joint calibration set of size $K \cdot N$. To calibrate, the merged model is placed in PyTorch's native training mode ($\mathtt{train()}$), but all learnable parameters are frozen. The momentum of all BatchNorm layers is temporarily set to $1.0$.

We perform a single task-agnostic forward pass of $D_{\text{cal}}$ through the backbone:
\begin{equation}
    H_l = \Phi_l(D_{\text{cal}})
\end{equation}
where $\Phi_l$ represents the layers up to BatchNorm layer $l$. This forward pass triggers PyTorch's native BatchNorm mechanism, directly calculating the batch mean $\mu_l$ and batch variance $\sigma_l^2$ of the joint calibration set, and completely overwriting the merged running mean and variance statistics:
\begin{equation}
    \mu_{l,\text{running}} \leftarrow \mu_{l}
\end{equation}
\begin{equation}
    \sigma^2_{l,\text{running}} \leftarrow \sigma^2_l
\end{equation}
This beautifully restores correct activation scales and mean shifts throughout the backbone task-agnostically. Since BatchNorm layers are positioned prior to ReLU activations in standard architectures, this native calibration naturally preserves ReLU sparsity, as ReLU continues to act as a clean sparsity filter on the calibrated representation space.

\subsubsection{Step 2: N-TAAC-Guided Head Adaptation}
Once the backbone statistics are updated, we freeze the backbone (including the calibrated BatchNorm statistics) and optimize only the task classification heads $h_k$ on $D_{\text{cal}}^{(k)}$. This is performed under two settings:
\begin{itemize}
    \item \textbf{Supervised SFT (SPJA-SFT):} We optimize the head $h_k$ using cross-entropy loss with labeled samples:
    \begin{equation}
        \label{eq:sft}
        \mathcal{L}_{\text{CE}} = -\frac{1}{N}\sum_{(x, y) \in D_{\text{cal}}^{(k)}} \log \text{softmax}(h_k(\hat{\Phi}(x)))_y
    \end{equation}
    where $\hat{\Phi}(x)$ represents features extracted from the N-TAAC-calibrated backbone.
    \item \textbf{Unsupervised TTA (SPJA-TTA):} If labels are unavailable, we adapt the head $h_k$ via soft KL-divergence distillation from the frozen expert teacher model $f_k^{\text{orig}}$:
    \begin{equation}
        \label{eq:tta}
        \begin{aligned}
        \mathcal{L}_{\text{KL}} = \frac{1}{N}\sum_{x \in D_{\text{cal}}^{(k)}} \text{KL}( &\text{softmax}(h_k(\hat{\Phi}(x))) \\
        &\parallel \text{softmax}(f_k^{\text{orig}}(x)) )
        \end{aligned}
    \end{equation}
\end{itemize}
By optimizing the head directly in the joint calibrated representation space of N-TAAC, the decision boundaries are perfectly aligned with the restored feature distributions, resolving task conflicts and localized head misalignment comprehensively without requiring any custom hooks during inference.

The complete SPJA procedure is summarized in Algorithm~\ref{alg:spja}.

\begin{algorithm*}[t]
  \caption{Sparsity-Preserving Joint Alignment (SPJA)}
  \label{alg:spja}
  \begin{algorithmic}[1]
    \STATE {\bfseries Input:} Pre-trained base model $\Theta_{\text{pre}}$, expert models $\{\Theta_1, \dots, \Theta_K\}$ with backbones $\theta_k$ and heads $h_k$; calibration datasets $\{D_{\text{cal}}^{(1)}, \dots, D_{\text{cal}}^{(K)}\}$ with size $N$ per task; scaling coefficient $\lambda$; learning rate $\eta$; training epochs $E$.
    \STATE {\bfseries Output:} Calibrated merged backbone $\theta_{\text{merged}}$, adapted task heads $\{h_1, \dots, h_K\}$.
    \STATE \COMMENT{\textbf{Step 1: Backbone Merging}}
    \IF{using Weight Averaging (WA)}
      \STATE $\theta_{\text{merged}} \leftarrow \frac{1}{K} \sum_{k=1}^K \theta_k$
    \ELSIF{using Task Arithmetic (TA)}
      \STATE $\theta_{\text{merged}} \leftarrow \theta_{\text{pre}} + \lambda \sum_{k=1}^K (\theta_k - \theta_{\text{pre}})$
    \ENDIF
    \STATE \COMMENT{\textbf{Step 2: Joint Backbone Calibration (N-TAAC)}}
    \STATE Construct joint calibration set: $D_{\text{cal}} \leftarrow \cup_{k=1}^K D_{\text{cal}}^{(k)}$
    \STATE Freeze all learnable parameters in $\theta_{\text{merged}}$
    \STATE Set model to PyTorch training mode: $\theta_{\text{merged}}.\mathtt{train()}$
    \STATE Set BatchNorm statistic tracking momentum to $1.0$ for all layers
    \STATE Perform single task-agnostic forward pass: $\Phi(D_{\text{cal}}; \theta_{\text{merged}})$
    \STATE \COMMENT{BatchNorm statistics are directly updated in $\theta_{\text{merged}}$}
    \STATE Set model to evaluation mode: $\theta_{\text{merged}}.\mathtt{eval()}$
    \STATE \COMMENT{\textbf{Step 3: N-TAAC-Guided Head Adaptation}}
    \FOR{each task $k \in \{1, \dots, K\}$}
      \STATE Extract calibrated features: $\hat{\Phi}(D_{\text{cal}}^{(k)}) \leftarrow \Phi(D_{\text{cal}}^{(k)}; \theta_{\text{merged}})$
      \FOR{epoch $e = 1$ {\bfseries to} $E$}
        \IF{Supervised SFT (SPJA-SFT)}
          \STATE Compute cross-entropy loss $\mathcal{L}_{\text{CE}}$ as in Eq.~\ref{eq:sft}
        \ELSIF{Unsupervised TTA (SPJA-TTA)}
          \STATE Compute teacher distillation loss $\mathcal{L}_{\text{KL}}$ as in Eq.~\ref{eq:tta}
        \ENDIF
        \STATE Update head: $h_k \leftarrow h_k - \eta \nabla_{h_k} \mathcal{L}$
      \ENDFOR
    \ENDFOR
    \STATE {\bfseries return} $\theta_{\text{merged}}, \{h_1, \dots, h_K\}$
  \end{algorithmic}
\end{algorithm*}

\section{Experimental Setup}
We employ the standard multi-task vision benchmark comprising \textbf{MNIST} \cite{lecun1998gradient}, \textbf{Fashion-MNIST} \cite{xiao2017fashion}, and \textbf{CIFAR-10} \cite{krizhevsky2009learning} to test our hypothesis.

\subsection{Backbone Architecture and Training}
Following prior work, we employ a standard ResNet-18 \cite{he2016deep} pre-trained on ImageNet-1k \cite{deng2009imagenet}. Grayscale datasets (MNIST, Fashion-MNIST) are converted to 3-channel RGB and resized to $32 \times 32$ to share the backbone. The ResNet-18 classification layer is replaced by $K=3$ task-specific heads. Each expert is fine-tuned independently for 5 epochs using AdamW \cite{loshchilov2017decoupled} with a learning rate of $5\times 10^{-4}$, weight decay of $10^{-4}$, batch size of 128, Cosine Annealing, and standard data augmentations (random crop and horizontal flip).

\subsection{Evaluation Protocols and Baselines}
We compare our framework, \textbf{Sparsity-Preserving Joint Alignment (SPJA)}, in supervised (SPJA-SFT) and unsupervised (SPJA-TTA) modes against five baselines under identical data splits and budgets:
\begin{enumerate}
    \item \textbf{NONE:} Direct backbone merging using WA or TA, evaluated with original heads.
    \item \textbf{N-TAAC:} Updates BatchNorm statistics on a joint calibration set in a single forward pass in native \texttt{train()} mode with frozen weights.
    \item \textbf{LSC:} Scales activations globally per layer using a single positive standard deviation ratio.
    \item \textbf{Head-SFT:} Freezes the backbone and optimizes only task heads on labeled calibration subsets.
    \item \textbf{Head-TTA:} Freezes the backbone and adapts only task heads using soft teacher self-distillation.
\end{enumerate}

For all calibration and adaptation methods, we use a default calibration size of $N=128$ per task, and optimize heads for 15 epochs using AdamW with a learning rate of $10^{-3}$ and batch size equal to the calibration size.

\section{Results and Empirical Analysis}
We report the classification accuracy (\%) over each test set and compute multi-task performance via average accuracy across all three tasks.

\subsection{Weight Merging Results}
Our main results under Weight Averaging (WA) are summarized in Table \ref{tab:wa_results}.
\begin{table}[h]
\caption{Multi-task model merging accuracies (\%) on ResNet-18 vision benchmark under Weight Averaging (WA) with calibration size $N=128$ and random seed $42$. Best merging results are in bold.}
\label{tab:wa_results}
\begin{center}
\begin{small}
\begin{sc}
\resizebox{\columnwidth}{!}{%
\begin{tabular}{lcccc}
\toprule
Method / Model & MNIST & Fashion-MNIST & CIFAR-10 & Average \\
\midrule
Individual Experts (Upper Bound) & 99.06 & 92.20 & 85.11 & 92.12 \\
\midrule
WA Baseline (NONE) & 43.81 & 15.32 & 28.53 & 29.22 \\
N-TAAC & 88.54 & 76.62 & 52.61 & 72.59 \\
LSC & 11.35 & 13.98 & 10.48 & 11.94 \\
Head-SFT & 75.17 & 59.56 & 39.85 & 58.19 \\
Head-TTA & 74.30 & 60.73 & 37.30 & 57.44 \\
\midrule
SPJA-SFT (Ours) & 92.07 & 80.10 & 56.71 & 76.29 \\
SPJA-TTA (Ours) & \textbf{92.41} & \textbf{80.22} & \textbf{57.49} & \textbf{76.71} \\
\bottomrule
\end{tabular}%
}
\end{sc}
\end{small}
\end{center}
\vskip -0.1in
\end{table}
 Our findings demonstrate a remarkable synergy in Sparsity-Preserving Joint Alignment. While baseline Weight Averaging (NONE) collapses catastrophically, and individual methods like LSC or Head Adaptation achieve moderate recovery, SPJA-SFT and SPJA-TTA achieve spectacular performance gains, validating our core hypothesis.

\begin{figure*}[t]
\centering
\begin{subfigure}{0.32\textwidth}
  \centering
  \includegraphics[width=\linewidth]{fig_lambda_sweep.pdf}
  \caption{Task Arithmetic $\lambda$ Sweep}
  \label{fig:lambda_sweep}
\end{subfigure}
\hfill
\begin{subfigure}{0.32\textwidth}
  \centering
  \includegraphics[width=\linewidth]{fig_convergence.pdf}
  \caption{Convergence Analysis}
  \label{fig:convergence}
\end{subfigure}
\hfill
\begin{subfigure}{0.32\textwidth}
  \centering
  \includegraphics[width=\linewidth]{fig_cal_strategy.pdf}
  \caption{Calibration Selection Strategy}
  \label{fig:cal_strategy}
\end{subfigure}
\caption{Comprehensive empirical evaluation of Sparsity-Preserving Joint Alignment (SPJA) showing (a) robust performance recovery across all Task Arithmetic scaling factors, (b) rapid and stable convergence over head-adaptation epochs, and (c) superior data efficiency of Class-Balanced calibration selection.}
\label{fig:empirical_evaluation}
\end{figure*}

\subsection{Task Arithmetic and Hyperparameter Sweeps}
We investigate the performance of our proposed method under Task Arithmetic (TA) as a function of the scaling coefficient $\lambda \in \{0.1, 0.2, 0.3, 0.4\}$. Table \ref{tab:ta_results} presents the average multi-task accuracies, and Figure~\ref{fig:lambda_sweep} illustrates this performance trend.
\begin{table}[h]
\caption{Multi-task model merging average accuracies (\%) under Task Arithmetic (TA) across varying scaling coefficients $\lambda$ on ResNet-18.}
\label{tab:ta_results}
\begin{center}
\begin{small}
\begin{sc}
\resizebox{\columnwidth}{!}{%
\begin{tabular}{lcccc}
\toprule
Method / Model & $\lambda = 0.1$ & $\lambda = 0.2$ & $\lambda = 0.3$ & $\lambda = 0.4$ \\
\midrule
TA Baseline (NONE) & 28.80 & 40.36 & 37.95 & 9.93 \\
N-TAAC & 44.69 & 64.32 & 71.44 & 73.94 \\
LSC & 10.38 & 12.54 & 12.23 & 9.93 \\
Head-SFT & 48.37 & 57.30 & 61.67 & 9.93 \\
Head-TTA & 48.50 & 57.95 & 61.88 & 9.93 \\
\midrule
SPJA-SFT (Ours) & 58.52 & 69.61 & 75.35 & 77.44 \\
SPJA-TTA (Ours) & \textbf{58.77} & \textbf{70.15} & \textbf{75.80} & \textbf{77.62} \\
\bottomrule
\end{tabular}%
}
\end{sc}
\end{small}
\end{center}
\vskip -0.1in
\end{table}
 Once again, SPJA consistently dominates both the uncalibrated baseline and individual methods across all scaling regimes.

\begin{table*}[t]
\begin{minipage}{0.48\textwidth}
\caption{Multi-task average accuracy (\%) of SPJA under varying Calibration Selection Strategies across sizes $N$ under Weight Averaging (WA).}
\label{tab:cal_strategy}
\begin{center}
\begin{small}
\begin{sc}
\resizebox{\columnwidth}{!}{%
\begin{tabular}{llccc}
\toprule
Method & Strategy & $N=32$ & $N=128$ & $N=512$ \\
\midrule
SPJA-SFT & Sequential & 71.51 & 76.29 & 78.26 \\
& Random & 69.38 & 76.57 & 78.18 \\
& Balanced & \textbf{73.30} & \textbf{76.79} & \textbf{78.42} \\
\midrule
SPJA-TTA & Sequential & 73.41 & 76.71 & 78.54 \\
& Random & 72.95 & 76.84 & 78.50 \\
& Balanced & \textbf{74.70} & \textbf{77.26} & \textbf{78.57} \\
\bottomrule
\end{tabular}%
}
\end{sc}
\end{small}
\end{center}
\end{minipage}
\hfill
\begin{minipage}{0.48\textwidth}
\caption{Multi-task average accuracy (\%) under Calibration Task Skew benchmarks on ResNet-18. We compare balanced data against three severe task imbalance configurations (represented as calibration dataset sizes).}
\label{tab:calibration_skew}
\begin{center}
\begin{small}
\begin{sc}
\resizebox{\columnwidth}{!}{%
\begin{tabular}{lcccc}
\toprule
Method / Skew & Balanced & MNIST-Heavy & Fashion-Heavy & CIFAR-Heavy \\
& (128, 128, 128) & (128, 13, 13) & (13, 128, 13) & (13, 13, 128) \\
\midrule
NONE Baseline & 29.22 & 29.22 & 29.22 & 29.22 \\
N-TAAC alone & 72.59 & 67.60 & 68.71 & 63.61 \\
LSC alone & 11.94 & 11.89 & 11.94 & 11.89 \\
\midrule
Head-SFT alone & 58.19 & 46.11 & 46.21 & 39.27 \\
Head-TTA alone & 57.44 & 44.52 & 46.10 & 37.24 \\
\midrule
SPJA-SFT (Ours) & 76.29 & 63.57 & 65.13 & 59.53 \\
\textbf{SPJA-TTA (Ours)} & \textbf{76.71} & \textbf{67.21} & \textbf{65.44} & \textbf{62.11} \\
\bottomrule
\end{tabular}%
}
\end{sc}
\end{small}
\end{center}
\end{minipage}
\vskip -0.1in
\end{table*}

\subsection{Sample Complexity Ablation Study}
To evaluate data efficiency, we sweep the calibration dataset size $N \in \{32, 128, 512\}$ samples per task. As shown in Table \ref{tab:sample_complexity}, SPJA remains highly robust even in extremely low-data regimes ($N = 32$), confirming its practical utility for real-world deployments where data is scarce.

\subsection{Robustness and Statistical Significance}
To ensure the reproducibility of our findings, we evaluate WA performance across three independent random seeds (42, 100, and 2026). As summarized in Table \ref{tab:seed_robustness}, the performance of SPJA is exceptionally stable, exhibiting a standard deviation of less than 0.5\% across all seeds, confirming that the benefits of joint alignment are highly statistically significant.

\begin{table*}[t]
\begin{minipage}{0.48\textwidth}
\caption{Multi-task average accuracy (\%) as a function of calibration set size per task $N$ under Weight Averaging (WA).}
\label{tab:sample_complexity}
\begin{center}
\begin{small}
\begin{sc}
\resizebox{\columnwidth}{!}{%
\begin{tabular}{lccc}
\toprule
Method / Model & $N=32$ & $N=128$ & $N=512$ \\
\midrule
N-TAAC & 71.08 & 72.59 & 72.82 \\
Head-SFT & 43.94 & 58.19 & 64.84 \\
Head-TTA & 44.64 & 57.44 & 65.30 \\
\midrule
SPJA-SFT (Ours) & 71.51 & 76.29 & 78.26 \\
SPJA-TTA (Ours) & \textbf{72.52} & \textbf{76.71} & \textbf{78.54} \\
\bottomrule
\end{tabular}%
}
\end{sc}
\end{small}
\end{center}
\end{minipage}
\hfill
\begin{minipage}{0.48\textwidth}
\caption{Multi-task average accuracy (\%) across three independent random seeds under Weight Averaging (WA) with calibration size $N=128$.}
\label{tab:seed_robustness}
\begin{center}
\begin{small}
\begin{sc}
\resizebox{\columnwidth}{!}{%
\begin{tabular}{lcccc}
\toprule
Method / Model & Seed 42 & Seed 100 & Seed 2026 & Mean $\pm$ Std \\
\midrule
Head-SFT & 58.19 & 58.27 & 57.67 & 58.04 $\pm$ 0.27 \\
Head-TTA & 57.44 & 59.22 & 57.68 & 58.11 $\pm$ 0.79 \\
\midrule
SPJA-SFT (Ours) & 76.29 & 76.28 & 75.97 & 76.18 $\pm$ 0.15 \\
SPJA-TTA (Ours) & \textbf{76.71} & \textbf{77.15} & \textbf{76.84} & \textbf{76.90} $\pm$ 0.19 \\
\bottomrule
\end{tabular}%
}
\end{sc}
\end{small}
\end{center}
\end{minipage}
\vskip -0.1in
\end{table*}

\subsection{Hyperparameter Sensitivity and Robustness}
To fully validate the stability of Sparsity-Preserving Joint Alignment (SPJA), we conduct exhaustive hyperparameter sweeps over the classification head optimization process. True to an empirically driven evaluation, we analyze the sensitivity of SPJA under: (i) varying learning rates $\eta \in \{10^{-4}$, $5\times 10^{-4}$, $10^{-3}$, $2\times 10^{-3}\}$ at a fixed epoch budget of 15, and (ii) varying training epochs $E \in \{5, 10, 15, 20\}$ at a fixed learning rate of $10^{-3}$. These sweeps are evaluated under Weight Averaging (WA) with a calibration size of $N=128$.

The results of the learning rate sweep are summarized in Table \ref{tab:lr_sweep}. We find that SPJA exhibits remarkable stability across an entire order of magnitude of learning rates. Under all tested configurations, SPJA-TTA consistently outperforms both the uncalibrated baseline (29.22\%) and individual methods, remaining highly robust to optimization hyperparameter choices.

Table \ref{tab:epoch_ablation} presents the convergence analysis across varying head-adaptation epochs, which is visually represented in Figure~\ref{fig:convergence}. Remarkably, SPJA recovers substantial performance ($>74.7\%$) in just 5 epochs (equivalent to 5 gradient steps on 128 samples per task), highlighting its incredible computational and sample efficiency. The performance peaks around 15 epochs and stabilizes without showing any signs of overfitting by epoch 20.

\begin{table*}[t]
\begin{minipage}{0.48\textwidth}
\caption{Multi-task average accuracy (\%) of SPJA as a function of the head-adaptation learning rate $\eta$ under Weight Averaging (WA) with $N=128$.}
\label{tab:lr_sweep}
\begin{center}
\begin{small}
\begin{sc}
\resizebox{\columnwidth}{!}{%
\begin{tabular}{lcc}
\toprule
Learning Rate $\eta$ & SPJA-SFT (Ours) & SPJA-TTA (Ours) \\
\midrule
$1 \times 10^{-4}$ & 74.51 & 74.58 \\
$5 \times 10^{-4}$ & 75.23 & 76.04 \\
$1 \times 10^{-3}$ & 76.13 & \textbf{76.78} \\
$2 \times 10^{-3}$ & 74.46 & 74.62 \\
\bottomrule
\end{tabular}%
}
\end{sc}
\end{small}
\end{center}
\end{minipage}
\hfill
\begin{minipage}{0.48\textwidth}
\caption{Multi-task average accuracy (\%) of SPJA as a function of head-adaptation epochs $E$ under Weight Averaging (WA) with $N=128$.}
\label{tab:epoch_ablation}
\begin{center}
\begin{small}
\begin{sc}
\resizebox{\columnwidth}{!}{%
\begin{tabular}{lcc}
\toprule
Epochs $E$ & SPJA-SFT (Ours) & SPJA-TTA (Ours) \\
\midrule
5 & 74.71 & 74.80 \\
10 & 75.64 & 76.09 \\
15 & 76.13 & \textbf{76.78} \\
20 & 75.73 & 76.66 \\
\bottomrule
\end{tabular}%
}
\end{sc}
\end{small}
\end{center}
\end{minipage}
\vskip -0.1in
\end{table*}

\subsection{Parameter Efficiency and Adaptation Depth}
To further deconstruct the success of Head-only adaptation in Sparsity-Preserving Joint Alignment, we conduct a systematic ablation study on the depth of the trainable layers during joint alignment. Specifically, we compare three levels of adaptation depth $\mathcal{D}$:
\begin{itemize}
    \item \textbf{Heads Only ($\mathcal{D}=\text{Heads}$):} Adapting only the task-specific classification heads (15k parameters), which is the default SPJA configuration.
    \item \textbf{Layer 4 + Heads ($\mathcal{D}=\text{Layer4}$):} Adapting the final residual block (\texttt{layer4}) of the backbone plus the classification heads (4.7M parameters).
    \item \textbf{Full Backbone + Heads ($\mathcal{D}=\text{Full}$):} Fine-tuning the entire network (11.2M parameters) sequentially on task calibration sets.
\end{itemize}

The results of this ablation study are presented in Table~\ref{tab:adaptation_depth}.

\begin{table}[h]
\caption{Adaptation depth ablation study under Weight Averaging (WA) with calibration size $N=128$. We compare adapting only classification Heads, Layer 4 + Heads, and the Full backbone + Heads.}
\label{tab:adaptation_depth}
\begin{center}
\begin{small}
\begin{sc}
\resizebox{\columnwidth}{!}{%
\begin{tabular}{lcccc}
\toprule
Method / Depth & MNIST & Fashion-MNIST & CIFAR-10 & Average \\
\midrule
\multicolumn{5}{l}{\textbf{Adaptation Depth: Heads}} \\
\midrule
Head-SFT & 75.17 & 59.56 & 39.85 & 58.19 \\
Head-TTA & 74.30 & 60.73 & 37.30 & 57.44 \\
SPJA-SFT (Ours) & 92.07 & 80.10 & 56.71 & 76.29 \\
\textbf{SPJA-TTA (Ours)} & \textbf{92.41} & \textbf{80.22} & \textbf{57.49} & \textbf{76.71} \\
\midrule
\multicolumn{5}{l}{\textbf{Adaptation Depth: Layer4}} \\
\midrule
Head-SFT & 45.99 & 69.98 & 49.92 & 55.30 \\
Head-TTA & 54.59 & 71.07 & 56.55 & 60.74 \\
SPJA-SFT (Ours) & 49.15 & 61.49 & 52.65 & 54.43 \\
\textbf{SPJA-TTA (Ours)} & 53.97 & 72.27 & 51.58 & 59.27 \\
\midrule
\multicolumn{5}{l}{\textbf{Adaptation Depth: Full}} \\
\midrule
Head-SFT & 20.78 & 50.86 & 40.44 & 37.36 \\
Head-TTA & 17.81 & 57.80 & 41.95 & 39.19 \\
SPJA-SFT (Ours) & 23.17 & 54.92 & 38.02 & 38.70 \\
\textbf{SPJA-TTA (Ours)} & 20.86 & 45.36 & 42.10 & 36.11 \\
\bottomrule
\end{tabular}%
}
\end{sc}
\end{small}
\end{center}
\vskip -0.1in
\end{table}

We observe a severe and catastrophic collapse in performance when adapting shared backbone layers ($\mathcal{D}=\text{Layer4}$ or $\mathcal{D}=\text{Full}$). Because the calibration sets are processed sequentially (Task 0, then Task 1, then Task 2), updating shared backbone weights task-specifically causes extreme catastrophic interference. During sequential fine-tuning, the updates on later tasks (e.g., CIFAR-10) entirely destroy the feature representations optimized or restored for earlier tasks (e.g., MNIST). Furthermore, with only $N=128$ samples per task, updating millions of parameters leads to severe overfitting.

In stark contrast, our default heads-only configuration ($\mathcal{D}=\text{Heads}$) completely isolates task-specific updates, ensuring zero sequential interference while requiring 750$\times$ fewer parameter updates. This demonstrates that head-only adaptation is not only parameter-efficient, but is also a structural requirement for multi-task joint alignment under data scarcity.

\subsection{Calibration Selection Strategies}
To evaluate the impact of calibration data composition on SPJA, we compare three selection strategies under Weight Averaging (WA) across $N \in \{32, 128, 512\}$: (i) \textbf{Sequential} (taking first $N$ training samples), (ii) \textbf{Random} (unstructured random sampling), and (iii) \textbf{Balanced} (uniform class-balanced sampling).

The multi-task average accuracies are summarized in Table~\ref{tab:cal_strategy} and visually analyzed in Figure~\ref{fig:cal_strategy}. We observe that Class-Balanced selection systematically outperforms both sequential and random sampling. The benefit of class balance is particularly prominent under severe data scarcity ($N = 32$), where Balanced SPJA-TTA achieves \textbf{74.70}\%, an absolute improvement of +1.29\% and +1.75\% over sequential and random sampling respectively. As $N$ increases to 512, the performance differences converge as expected. This highlights that guaranteeing class balance is a critical and highly data-efficient strategy for aligning decision boundaries.

\subsection{Robustness to Calibration Task Skew}
In practical deployments, multi-task merged models are often subjected to highly asymmetric and skewed data distributions, where data may be scarce or unavailable for certain tasks. To evaluate the resilience of our framework to severe task imbalance and data scarcity during alignment, we simulate a \textit{Calibration Task Skew} benchmark. 

We utilize a default calibration scale of $N=128$, and introduce a skewness factor $r_k \in \{1.0, 0.1\}$ for each task $k \in \{1, 2, 3\}$. A multiplier of $r_k=0.1$ reduces task $k$'s calibration dataset size to just $N_k = 13$ samples, while $r_k=1.0$ maintains the full $128$ samples. We evaluate under four asymmetric configurations: Balanced ($128, 128, 128$ samples); MNIST-Heavy ($128, 13, 13$ samples); Fashion-Heavy ($13, 128, 13$ samples); and CIFAR-Heavy ($13, 13, 128$ samples). Both joint backbone calibration (N-TAAC) and task-specific classification head adaptation are performed using these skewed calibration sets.

The results are summarized in Table~\ref{tab:calibration_skew}. We observe that existing post-hoc head-adaptation baselines (Head-SFT and Head-TTA) collapse catastrophically under task skew, struggling to align decision boundaries on a distorted, variance-collapsed backbone when task-specific data is extremely limited ($N_k=13$). For instance, under CIFAR-Heavy skew, Head-TTA alone drops to an average accuracy of just \textbf{37.24}\%.

In stark contrast, SPJA-TTA exhibits exceptional robustness, achieving \textbf{62.11}\% average accuracy under the same severe CIFAR-Heavy skew—a massive \textbf{+24.87}\% absolute performance gain over the baseline. Under MNIST-Heavy and Fashion-Heavy skews, SPJA-TTA achieves similar spectacular gains (+22.69\% and +19.34\% absolute improvements, respectively). 

This outstanding robustness stems from SPJA's decoupled architecture: joint backbone calibration (N-TAAC) successfully stabilizes the representational flow globally and task-agnostically, leaving a clean, healthy, and high-signal representation space. This enables even extremely low-data head optimization (TTA via soft teacher-distillation on just 13 samples) to easily find stable, high-generalizing decision boundaries, completely avoiding overfitting and confirming SPJA as a highly practical, robust paradigm.

\section{Discussion and Analysis}
The success of SPJA deconstructs the merging process, demonstrating that representation calibration and decision-boundary alignment are deeply complementary rather than competing paradigms.

While parameter averaging collapses backbone activation variance, standalone calibration cannot resolve localized head conflicts in the final classification layers. Conversely, standalone head adaptation is severely limited by optimization on a distorted, variance-collapsed feature space.

SPJA simultaneously solves both issues: N-TAAC provides a stable, task-agnostic, and variance-restored backbone representation space, enabling efficient head adaptation to align task decision boundaries with minimal data while maintaining strict mathematical consistency during training and inference.

\subsection{Limitations and Scope}
While Sparsity-Preserving Joint Alignment (SPJA) offers exceptional performance gains, we recognize certain limitations and areas for future work:
\begin{enumerate}
    \item \textbf{Calibration Data Dependency:} SPJA requires a small calibration dataset ($N \ge 32$ samples per task) to calculate backbone statistics and optimize heads. While extremely data-efficient and practical for most deployment settings, SPJA cannot be applied in purely zero-shot merging scenarios where absolutely no calibration samples are available.
    \item \textbf{BatchNorm Architecture Focus:} Our current implementation of N-TAAC leverages BatchNorm layers to capture and update the joint activation distribution globally. In architectures that utilize LayerNorm or RMSNorm (such as Vision Transformers or Large Language Models), standard layer normalization does not track running statistics for inference. Extending N-TAAC to Transformer backbones requires either introducing light-weight scale adaptation parameters or applying a scaling-only normalization alignment scheme.
    \item \textbf{Heads Parameter Scaling:} Although classification head adaptation updates $750\times$ fewer parameters than full fine-tuning, the number of task heads grows linearly with the number of merged tasks $K$. In extremely large-scale multi-task setups (e.g., hundreds of tasks), managing a large collection of classification heads may introduce minor memory overhead, though it remains vastly cheaper than deploying individual backbone models.
\end{enumerate}

\section{Conclusion and Future Work}
In this work, we deconstructed the paradigms of representation calibration and head-only adaptation in multi-task model merging. We introduced Sparsity-Preserving Joint Alignment (SPJA), a simple, stable, and highly effective framework that combines Native Task-Agnostic Activation Calibration (N-TAAC) with decision-boundary head adaptation. Through exhaustive, multi-seed empirical evaluations, we demonstrated that SPJA consistently outperforms existing baselines, establishing a new state-of-the-art standard. Future work includes extending SPJA to transformer architectures and exploring its performance under severe domain shifts.

\section*{Impact Statement}
This paper presents work whose goal is to advance the field of Machine Learning. There are many potential societal consequences of our work, none of which we feel must be specifically highlighted here.

\clearpage
\bibliography{submission}
\bibliographystyle{icml2026}

\end{document}

% This document was modified from the file originally made available by
% Pat Langley and Andrea Danyluk for ICML-2K. This version was created
% by Iain Murray in 2018, and modified by Alexandre Bouchard in
% 2019 and 2021 and by Csaba Szepesvari, Gang Niu and Sivan Sabato in 2022.
% Modified again in 2023 and 2024 by Sivan Sabato and Jonathan Scarlett.
% Previous contributors include Dan Roy, Lise Getoor and Tobias
% Scheffer, which was slightly modified from the 2010 version by
% Thorsten Joachims & Johannes Fuernkranz, slightly modified from the
% 2009 version by Kiri Wagstaff and Sam Roweis's 2008 version, which is
% slightly modified from Prasad Tadepalli's 2007 version which is a
% lightly changed version of the previous year's version by Andrew
% Moore, which was in turn edited from those of Kristian Kersting and
% Codrina Lauth. Alex Smola contributed to the algorithmic style files.
